\textbf{Commun à tous les candidats}

\medskip











\begin{enumerate}
\item 

\[f(x) = \left(x^2 - 2x - 1\right)\e^x.\]






On a $f(x) = x^2\e^x - 2x\e^x - \e^x$.

\starredbullet~ On sait que $\displaystyle\lim_{x \to - \infty} \e^x = 0$, puis que 

\starredbullet~$\displaystyle\lim_{x \to - \infty} 2x\e^x = 0$

\starredbullet~$\displaystyle\lim_{x \to - \infty} x^2\e^x = 0$, d'où par somme de limites : $\displaystyle\lim_{x \to - \infty} f(x) = 0$. Réponse C.
\item  On considère la fonction $f$ définie sur $\R$ par $f(x) = \dfrac{3}{5 + \e^x}$.









\starredbullet~On a $\displaystyle\lim_{x \to - \infty} f(x) = \dfrac{3}{5}$ : la droite d'équation $y = \dfrac{3}{5}$ est asymptote horizontale au voisinage de moins l'infini ;

\starredbullet~On a $\displaystyle\lim_{x \to + \infty} f(x) = 0$ : l'axe des abscisses est asymptote horizontale au voisinage de plus l'infini. Réponse C.
\item  
















On voit sur la figure que $f''(- 3) = f''(2) = f''(5) = 0$ : la dérivée seconde s'annule trois fois donc la fonction $f$ admet trois points d'inflexion. Réponse B.








$n^2 - 17n + 20 = \left(n - \frac{17}{2} \right)^2 - \left(\frac{17}{2}\right)^2 + 20 = \left(n - \frac{17}{2} \right)^2 - \frac{289}{4} + \frac{80}{4} = \left(n - \frac{17}{2} \right)^2 - \frac{209}{4} = \left(n - \frac{17 - \sqrt{209}}{2} \right)\left(n - \frac{17 + \sqrt{209}}{2}\right).$

On a donc quel que soit $n$, \, $u_n \geqslant - \frac{209}{4}$ : la suite est donc minorée. Réponse A.
\item 
























Réponse A.
\end{enumerate}

\bigskip
